\chapter{Fundamentação Teórica}

% Guardrails do capitulo (fonte: docs/living-paper/chapter_structure_freeze.md)
% Objetivo: fundamentar teoricamente previsao financeira com TFT, incerteza quantilica e avaliacao reprodutivel.
% Escopo obrigatorio: forecasting financeiro, deep learning temporal, TFT, OOS e reproducibilidade.
% Fora de escopo: revisao centrada em NLP/CNN para classificacao de sentimento como problema principal.
% Fonte primaria deste capitulo: docs/07_reports/living-paper/10_problem_and_related_work.md (v2 prosa).

\section{Previsão de séries temporais financeiras}
A previsão de retornos em mercados financeiros é historicamente tratada como um problema de fronteira entre teoria informacional e evidência empírica. Em sua segunda revisão da hipótese de mercados eficientes, \cite{FAMA1991} consolida a tese de que preços refletem rapidamente toda a informação pública disponível, de modo que retornos anormais sustentáveis dependem ou de informação privada, ou de exposição a fatores de risco não precificados, ou de fricções de mercado específicas. O corolário metodológico relevante para este trabalho é direto: previsibilidade estável de retornos lineares é a exceção, não a regra, em ativos líquidos.

Em contraponto, \cite{LO2004} formula a hipótese de mercados adaptativos, que recolhe os mesmos fatos estilizados sob uma lógica evolucionária: eficiência e ineficiência coexistem ao longo do tempo em função de competição entre estratégias, mudanças de regime e adaptação dos agentes. A consequência prática não é a refutação de Fama, mas o reconhecimento de que sinais preditivos -- quando existem -- tendem a ser contextuais, multivariados e temporalmente localizados, em vez de globalmente persistentes. Dois fenômenos amplamente documentados reforçam essa leitura: não estacionariedade de primeira e segunda ordem em retornos diários, e clustering de volatilidade.

Evidência empírica recente recalibra a expectativa de magnitude do sinal. \cite{GU2020}, em revisão sistemática de aprendizado de máquina aplicado a apreçamento de ativos, reportam $R^2$ fora da amostra tipicamente entre 0,1\% e 1\% para retornos diários de ações líquidas mesmo com modelos flexíveis e amplo conjunto de preditores. \cite{RAPACH2013}, em survey sobre previsão de retornos agregados, chegam a magnitudes compatíveis e discutem como pequenos $R^2$ OOS podem ainda ser economicamente relevantes quando combinados a alocação informada. Essa literatura ancora uma expectativa metodológica central para o presente trabalho: ganhos preditivos modestos são a norma, e qualquer protocolo que reporte ganhos sistematicamente altos deve ser tratado com suspeita de leakage ou data snooping antes da interpretação econômica.

A combinação de retornos pouco previsíveis, regimes não estacionários e $R^2$ OOS pequenos torna a robustez metodológica mais informativa do que o número pontual de erro: comparação pareada contra baselines explícitos, alinhamento temporal estrito, avaliação quantílica e correção para múltiplos testes deixam de ser ornamento acadêmico e passam a ser condição necessária para que o claim ``o modelo superou o baseline'' carregue significado. O escopo congelado deste trabalho (modelos single-asset, AAPL como piloto inicial, horizontes h+1 e h+7) reflete diretamente essa restrição: não se afirma generalização inter-ativos ou inter-períodos que a literatura não sustentaria.

\section{Modelagem com deep learning para séries temporais}
\cite{LIMZOHREN2021} sistematizam, em survey dedicado, o estado da arte em forecasting de séries temporais com aprendizado profundo, organizando famílias de arquiteturas (recorrentes, convolucionais, baseadas em atenção e híbridas) e papéis distintos que podem desempenhar (previsão pontual, multi-horizonte, probabilística, multivariada). Dois pontos da revisão são particularmente pertinentes a este trabalho. Primeiro, modelos profundos oferecem flexibilidade representacional para capturar não linearidades, interações entre covariáveis heterogêneas e dependências temporais de diferentes escalas, atributos diretamente relevantes para forecasting financeiro com indicadores técnicos, sentimento e fundamentos. Segundo, essa flexibilidade vem acompanhada de risco aumentado de overfitting e de sensibilidade a leakage temporal sutil, sobretudo quando o protocolo experimental não é desenhado com rigor.

A consequência metodológica reforça a discussão da seção anterior: a adoção de deep learning em forecasting financeiro não se justifica por premissa de superioridade universal, mas por aderência ao tipo de dado (multivariado, multi-escala, multi-horizonte) e por adequação a um protocolo de avaliação que controla explicitamente as fontes de inflação artificial de desempenho. Não basta declarar splits temporais; é preciso verificar leakage de features (scaler com fit em validação/teste, indicadores técnicos calculados com janela contendo dia atual, sentimento agregado fora da sessão de mercado etc.), leakage de seleção (hiperparâmetros escolhidos com o conjunto de teste) e leakage de alinhamento (intersecção temporal incompleta em comparações pareadas).

A escolha por uma arquitetura profunda neste trabalho é instrumental: serve para representar interações entre famílias de features e horizontes em um mesmo modelo, com mecanismos nativos de seleção de variáveis e atenção temporal. Ela não dispensa, mas torna mais crítica, a disciplina experimental de pré-registro, governança de coorte e alinhamento OOS exato.

\section{Temporal Fusion Transformer (TFT)}
O Temporal Fusion Transformer, proposto por \cite{LIMETAL2021TFT}, é um modelo de previsão multi-horizonte que combina componentes recorrentes e de atenção com mecanismos explícitos de seleção de variáveis e cabeça quantílica. Sua arquitetura organiza a inferência em estágios funcionais distintos: (i) processamento de variáveis estáticas via encoders específicos; (ii) Variable Selection Networks (VSN) que atribuem pesos dinâmicos por timestep e por família de feature, separando entradas estáticas, conhecidas no futuro (calendário, feriados planejados) e desconhecidas no futuro (preço, retorno, indicadores técnicos); (iii) blocos recorrentes LSTM no estágio encoder/decoder que carregam estado temporal entre passos; (iv) atenção multi-head que captura dependências de longo prazo entre períodos relevantes; e (v) cabeça multi-quantile que emite, por horizonte, múltiplos quantis condicionais do alvo em um único passo de inferência.

Três características justificam a escolha do TFT como modelo principal deste trabalho. Primeiro, o TFT integra covariáveis estruturadas heterogêneas em um único estimador, alinhado ao recorte de famílias PRICE/TECHNICAL/SENTIMENT/FUNDAMENTAL adotado no projeto. Segundo, ele é nativamente multi-horizonte, o que evita a necessidade de re-treinar modelos independentes por horizonte e simplifica a política de persistência adotada (uma linha por horizonte, mesmo checkpoint). Terceiro, sua cabeça quantílica produz, sem reajuste posterior, o objeto probabilístico exigido pelas hipóteses H1, H2a e H2b do trabalho: quantis condicionais por linha OOS, base para pinball loss e para métricas de calibração.

A interpretabilidade nativa do TFT, via VSN weights e atenção por timestep, é tratada neste trabalho como base do nível mais forte de evidência para análise de contribuição de features, conforme hierarquia metodológica adotada (global por horizonte $>$ agregada por regime $>$ local pontual). O próprio \cite{LIMETAL2021TFT} enfatiza, no entanto, que pesos de seleção indicam que variáveis o modelo usou sob um protocolo, ativo e período, e não quais variáveis causam o alvo -- distinção incorporada explicitamente como linguagem proibida no protocolo de explicabilidade.

Limites conhecidos. A capacidade representacional do TFT amplifica o risco discutido na seção anterior: hiperparâmetros, seed, particionamento e composição de features têm impacto material sobre a saída quantílica, incluindo patologias como quantile crossing e colapso de quantis ($p10 = p90$) observadas no analytics store deste projeto. Esses fenômenos motivam o pós-processamento por guardrail monotônico, que é parte integral do estimador entregue.

\section{Predição probabilística e quantílica}
A regressão quantílica, introduzida por \cite{KOENKERBASSETT1978}, é a base formal para estimação de quantis condicionais sem assumir forma paramétrica para a distribuição do erro. Em vez de minimizar o erro quadrático médio, estima-se cada quantil de interesse por minimização da pinball loss, função assimétrica que penaliza desvios em direções opostas com pesos distintos. A perda média de pinball entre quantis é a métrica primária deste trabalho para avaliação probabilística.

A justificativa para tratar pinball como métrica primária, e não apenas intervalos como PICP, vem do trabalho seminal de \cite{GNEITING2007} sobre scoring rules. Os autores estabelecem o conceito de regra estritamente própria, isto é, uma função de avaliação cuja minimização em esperança ocorre se e somente se a previsão reportada coincide com a distribuição verdadeira do alvo. Pinball loss é estritamente própria para o quantil correspondente, o que significa que reportar quantis honestos é estratégia dominante para o predictor; em contraste, métricas de cobertura agregada como PICP podem ser otimizadas com estratégias degeneradas (intervalos artificialmente largos) e portanto não são adequadas como função objetivo primária, ainda que sejam complementos necessários para diagnóstico de calibração.

A avaliação probabilística neste trabalho distingue formalmente duas formas de calibração: (a) calibração intervalar, medida por PICP do par ($p10$, $p90$) contra a cobertura nominal de 0,80; e (b) calibração marginal por quantil, medida pela taxa empírica $\textit{coverage}_{q\tau} = \mathrm{mean}(y_{\textit{true}} \le \textit{quantile}_{p\tau})$ contra o nominal $\tau$. Largura média (MPIW) é analisada conjuntamente com PICP: intervalo estreito com cobertura baixa indica modelo sistematicamente errado, não ``preciso''.

Patologia recorrente em regressores quantílicos treinados independentemente é o quantile crossing, situação em que $p10 > p50$ ou $p50 > p90$ em parte das linhas. \cite{CHERNOZHUKOV2010REARRANGEMENT} demonstram que o rearranjo monotônico da função quantil estimada produz estimador com perda esperada menor ou igual à do estimador original para qualquer perda convexa simétrica em torno da mediana, incluindo pinball, com igualdade somente quando o estimador original já é monotônico. Esse resultado teórico ancora a decisão deste trabalho de declarar a variante \textit{post-guardrail} (rearranjo via sort por linha) como primária para H1, H2a e H2b, com a variante \textit{raw} reportada em paralelo para auditar a magnitude do efeito.

Aplicação a risco. Quando os quantis previstos servem de base para métricas de risco financeiro, a monotonicidade não é apenas conveniência estatística: é pré-requisito de definição. \cite{JORION2007}, em referência clássica sobre Value-at-Risk, define VaR como quantil da distribuição de perda, exigindo que o estimador da função quantil seja monotônico para preservar a própria semântica de ``perda máxima ao nível $\alpha$''. \cite{ACERBI2002}, em discussão sobre coerência de medidas de risco, formalizam propriedades (monotonicidade, subaditividade, homogeneidade, invariância translacional) que Expected Shortfall satisfaz como medida coerente sob ordenação bem definida dos quantis. Sob crossing, nem VaR nem ES preservam essas propriedades; por isso, as métricas de risco deste trabalho usam exclusivamente a variante \textit{post-guardrail}.

\section{Avaliação de modelos em ambiente OOS e testes estatísticos}
Comparação entre modelos sob critério OOS exige, antes de qualquer teste estatístico, alinhamento temporal exato. Diferentes modelos podem ter inferências persistidas sob diferentes coberturas (por falha pontual de ingestão, por exigência de warmup distinto, por filtro de regime), e calcular médias de erro sobre amostras não idênticas confunde diferença de modelagem com diferença de cobertura. O regime formal adotado neste trabalho exige intersecção exata por (\textit{asset}, \textit{split}, \textit{horizon}, \textit{target\_timestamp}) antes de qualquer estatística pareada.

O teste de \cite{DIEBOLDMARIANO1995}, referência canônica para comparação de acurácia preditiva, opera sobre a série de perdas pareadas entre dois modelos e testa a hipótese nula de igualdade de perda esperada. Sua aplicação direta, no entanto, é conhecidamente conservadora ou anti-conservadora em amostras pequenas, dependendo da estrutura de autocorrelação das perdas. \cite{HARVEY1997} propõem a correção HLN para a estatística DM, baseada em ajuste de variância HAC e em distribuição $t$ com graus de liberdade adaptados ao tamanho efetivo da amostra; essa correção é o procedimento default usado pelos relatórios gold deste projeto para horizontes h+1 e h+7 sob amostra OOS limitada.

Comparações múltiplas entre o candidato TFT \textit{all-features} e os baselines admissíveis (random walk, AR(1), média histórica rolling, EWMA-vol, quantis empíricos históricos) inflam a taxa de erro tipo I quando reportadas sem correção. O ajuste adotado é Holm-Bonferroni, controle de \textit{Family-Wise Error Rate} por passos descendentes, persistido como \texttt{pvalue\_adj\_holm} nos artefatos gold. A regra de decisão para H2a (superioridade contra baseline primário) e H2b (superioridade contra todos os baselines simples) usa \texttt{pvalue\_adj\_holm $<$ 0,05} como critério formal de significância.

Quando se deseja, em vez de um vencedor pontual, identificar o conjunto de modelos estatisticamente indistinguíveis do melhor, recorre-se ao Model Confidence Set proposto por \cite{HANSEN2011MCS}. O procedimento controla FWER por construção e retorna o conjunto de configurações não rejeitadas como melhores ao nível $\alpha$, em vez de forçar um ranking único. Este trabalho usa MCS com $\alpha = 0{,}05$ no fluxo confirmatório da Fase B como leitura complementar a DM pareado: divergência entre as duas leituras (por exemplo, um modelo ganha em DM mas o MCS aceita vários modelos) é ela mesma evidência metodológica relevante.

A governança de coorte é parte do protocolo estatístico, não um detalhe operacional. Toda comparação confirmatória é restrita a runs com o mesmo \texttt{parent\_sweep\_id} (mesma coorte de hiperparâmetros e mesma fronteira temporal). Comparações globais sem restrição de coorte são usadas apenas como saúde de plataforma e não podem ser reportadas como evidência confirmatória.

\section{Reprodutibilidade experimental e rastreabilidade}
\cite{PINEAU2021REPRODUCIBILITY}, em relatório sobre o programa de reprodutibilidade da NeurIPS 2019, documentam que uma fração significativa de trabalhos publicados em conferências de ponta em aprendizado de máquina permanece difícil ou impossível de reproduzir mesmo com código aberto, em razão de omissões em (i) dados exatos de entrada, (ii) sementes e versões de biblioteca, (iii) hiperparâmetros, e (iv) procedimento exato de splits e seleção de modelo. Os autores propõem como mitigação um conjunto mínimo de artefatos obrigatórios, incluindo descrição explícita do pipeline, controle de aleatoriedade e disponibilidade de checkpoints. A consequência para trabalhos acadêmicos com pretensão de defensabilidade é clara: reprodutibilidade deixa de ser virtude operacional opcional e passa à condição necessária de validade.

Esse posicionamento ancora a política adotada neste trabalho. Reprodutibilidade aqui significa que decisões finais são reproduzíveis a partir de artefatos persistidos, sem necessidade de re-treinamento. Operacionalmente, isso implica três camadas. Primeiro, identidade lógica determinística por \textit{run}: o \texttt{run\_id} combina hashes de \textit{asset}, \texttt{feature\_set\_hash}, \texttt{trial\_number}, \textit{fold}, \textit{seed}, \texttt{model\_version}, \texttt{config\_signature}, \texttt{split\_signature} e \texttt{pipeline\_version} em um SHA-256 estável, de modo que o mesmo experimento sempre produz o mesmo identificador. Segundo, \textit{fingerprints} de dataset, configuração e \textit{split} asseguram que comparações pareadas estão sobre entradas idênticas. Terceiro, separação entre camadas imutáveis (fonte primária, \textit{silver}) e camadas reconstruíveis (\textit{gold}) permite que reanálises rodem sem retrain.

O pré-registro versionado é o componente final da política de reprodutibilidade. O pré-registro fixa, antes da execução da rodada confirmatória: candidato, baselines, métrica primária, contrato de variante quantílica (\texttt{primary\_quantile\_contract}), bandas de calibração, regra de desempate, threshold de relevância prática. O ato de versionar e commitar essa declaração antes de observar os resultados elimina graus de liberdade que, na ausência de pré-registro, geram p-hacking e seleção adversa de baselines.

A combinação de identidade determinística, \textit{fingerprints}, camadas imutáveis e pré-registro versionado permite que as conclusões do TCC sejam auditadas e reconstruídas em sessões futuras exclusivamente a partir do repositório em sua forma persistida. Essa propriedade não garante que as conclusões sejam corretas, mas garante que sejam falsificáveis no sentido operacional: qualquer revisor pode re-executar os procedimentos sobre os mesmos artefatos e chegar aos mesmos números, e qualquer divergência é por si só um achado.
